\section*{الفصل \ref{ch:b2:dispersion-wave-packets} --- التبدد والرزم الموجية}

\begin{solution}{exo:b2:dispersion-wave-packets:1}
(أ) $k = \qty{0.063}{rad/m}$، و $\omega = \sqrt{gk} = \qty{0.79}{rad/s}$، و $T = \qty{8.0}{s}$؛
و $v_\varphi = \qty{12.5}{m/s}$، و $v_g = \qty{6.3}{m/s}$. (ب) $\omega = \qty{0.45}{rad/s}$، و $k = \omega^2
/g = \qty{0.021}{rad/m}$، و $\lambda = \qty{310}{m}$؛ و $v_\varphi = gT/2\pi = \qty{22}{m/s}$، و $v_g =
\qty{11}{m/s}$؛ و $5 \times 10^6/11 = \qty{4.6e5}{s} = 5.3$ يوم. (ج) $\lambda = \qty{6}{m}$، و $v_g
= \qty{1.6}{m/s}$: فيلزمها خمسة أسابيع، والأمواج القصيرة تخمدها
\omterm{def:b2:viscous-flows:viscosity}{اللزوجة} وتكسرها الريح قبل ذلك بكثير.
\end{solution}

\begin{solution}{exo:b2:dispersion-wave-packets:2}
$\omega = \qty{1.26e8}{rad/s}$، و $\omega_c = \qty{6.3e7}{rad/s}$: ومنه $k = \sqrt{\omega^2 - \omega_c^2}/c =
\qty{0.36}{rad/m}$، و $v_\varphi = \qty{3.5e8}{m/s}$، و $v_g = c^2k/\omega = \qty{2.6e8}{m/s}$،
وجداؤهما $c^2$. وعند \qty{5}{MHz} يكون $\omega < \omega_c$: ومنه $\kappa = \sqrt{\omega_c^2 - \omega^2}/c =
\qty{0.18}{m^{-1}}$، والعمق \qty{5.5}{m}.
\end{solution}

\begin{solution}{exo:b2:dispersion-wave-packets:3}
الضربات عند \qty{4}{Hz}، ودور تغيّر العلوّ \qty{0.25}{s}. وأما الماء فتواتراه $f = 0.125$ و
\qty{0.119}{Hz}؛ والغلاف $\cos(\delta\omega t - \delta kx)$ مع $\delta k = \qty{2.9e-3}{rad/m}$
(أي ضربات تتباعد \qty{1.1}{km}) وسرعة $\delta\omega/\delta k = \qty{6.4}{m/s}$ --- وهي
\omterm{prop:b2:dispersion-wave-packets:group}{سرعة المجموعة} عند \qty{105}{m}، $\tfrac12\sqrt{g\lambda/2\pi} = \qty{6.4}{m/s}$.
\end{solution}

\begin{solution}{exo:b2:dispersion-wave-packets:4}
$\ln(b/a) = 1.20$: ومنه $\Lambda = \qty{0.24}{\micro H/m}$، و $\Gamma = \qty{106}{pF/m}$، و $c = 1/
\sqrt{\Lambda\Gamma} = \qty{2.0e8}{m/s}$، و $Z_c = \qty{48}{\ohm}$، و \qty{5.0}{ns/m}؛ و \qty{10}{ns}
$\leftrightarrow$ \qty{2.0}{m}. وبالهواء: $\Gamma = \qty{46}{pF/m}$، و $c = c_0$، و $Z_c = \qty{72}{\ohm}$.
\end{solution}

\begin{solution}{exo:b2:dispersion-wave-packets:5}
(أ) $\omega^2 = c^2k^2 + (EI/\mu)k^4$؛ و $v_\varphi = \sqrt{c^2 + (EI/\mu)k^2}$؛ و $v_g = (c^2k +
2(EI/\mu)k^3)/\omega$. (ب) $c^2 = \qty{1.12e5}{m^2/s^2}$، و $k = \qty{66}{rad/m}$، و $(EI/\mu)k^2
= \qty{7.0e3}{m^2/s^2}$: ومنه $v_\varphi/c = \sqrt{1.063} = 1.031$، أي $+3.1\%$ --- وهي $53$ سنتًا في
\cref{ch:b2:waves-on-strings}. (ج) $v_g > v_\varphi$: أي شاذّ؛ فتسبق التوافقيات
العالية وتصير مقدّمة القطفة زقزقةً هابطة قصيرة.
\end{solution}

\begin{solution}{exo:b2:dispersion-wave-packets:6}
(أ) $-\mu\omega^2 + \iu\alpha\omega = -Tk^2$: ومنه $\underline k^2 = (\omega^2/c^2)(1 - \iu\alpha/\mu\omega)$.
(ب) $\underline k \approx (\omega/c)(1 - \iu\alpha/2\mu\omega)$: أي $k' = \omega/c$، و $k'' = \alpha/2\mu c =
\alpha/2Z$؛ و $8.7k''$ \omterm{def:b2:sound-waves:decibel}{ديسيبل} في المتر. (ج) $c = \qty{100}{m/s}$، و $Z = \qty{0.5}{kg/s}$، و $\mu\omega = 3.1 \gg
0.02$: ومنه $k'' = \qty{0.02}{m^{-1}}$، و $\delta = \qty{50}{m}$، و \qty{0.17}{dB/m}. (د) لا في هذه
النهاية ($k' = \omega/c$)؛ وحين $\alpha \gtrsim \mu\omega$ لا يعود الجذر التربيعي
يتخطّى ويتوقف $k'$ على $\omega$ توقفًا غير خطي: فيصير متبدّدًا (وهي
نهاية الكابل الانتشارية).
\end{solution}

\begin{solution}{exo:b2:dispersion-wave-packets:7}
(أ) $v_\varphi = \sqrt{\gamma k/\rho}$، و $v_g = \tfrac32v_\varphi$: فتسبق المجموعة
الذرى. (ب) $v_\varphi^2 = g/k + \gamma k/\rho$ أصغري عند $k_m = \sqrt{\rho g/\gamma}$:
ومنه $\lambda_m = 2\pi\sqrt{\gamma/\rho g} = \qty{1.7}{cm}$، و $v_{\min} = (4g\gamma/\rho)^{1/4} =
\qty{23}{cm/s}$. (ج) والأثر نقش ساكن بالنسبة إلى
الزورق، مؤلَّف من أمواج تساوي سرعة طورها سرعته؛ ولا واحدة منها
أبطأ من \qty{23}{cm/s}. (د) والنظام الشعري، $v_g > v_\varphi$: فتولد الذرى
في مقدمة الحلقة وتموت في مؤخرتها.
\end{solution}

\begin{solution}{exo:b2:dispersion-wave-packets:8}
(أ) $k = \qty{0.126}{rad/m}$، و $\omega'' = -\tfrac14\sqrt{g/k^3} = \qty{-18}{m^2/s}$: ومنه $t_{\text{تمدد}}
= 2.5 \times 10^5/18 = \qty{1.4e4}{s}$، أي أربع ساعات. (ب) وعند \qty{10}{ps}: $t_{\text{تمدد}} =
4 \times 10^{-6}/0.16 = \qty{25}{\micro s}$، أي \qty{5}{km} عند \qty{2e8}{m/s}؛ وعند \qty{1}{ns}: $0.04/0.16 =
\qty{0.25}{s}$، أي \qty{50000}{km}. فالنبضات القصيرة (ومعدلات البتّات العالية) تتمدد في
كيلومترات: ومنه تستعمل الوصلات الطويلة ليفًا معوِّضًا أو معدلات معتدلة.
\end{solution}

\begin{solution}{exo:b2:dispersion-wave-packets:9}
(أ) $\partial_x^2v = R\Gamma\,\partial_tv$، و $D = 1/R\Gamma = \qty{1.1e12}{m^2/s}$. (ب) $\underline k^2
= -\iu R\Gamma\omega$: ومنه $k' = k'' = \sqrt{R\Gamma\omega/2}$؛ وعند \qty{1}{Hz}: $k = \qty{1.7e-6}{m^{-1}}$،
و $v_\varphi = \qty{3.7e6}{m/s}$، و $\delta = \qty{600}{km}$؛ وعند \qty{10}{Hz}: $v_\varphi = \qty{1.2e7}{m/s}$،
و $\delta = \qty{190}{km}$. (ج) $R\Gamma L^2 = 9 \times 10^{-13} \times 9 \times 10^{12} = \qty{8}{s}$
لكل نقطة: أي نحو كلمة في الدقيقة. (د) $\Lambda = R\Gamma/G = \qty{9}{H/km}$ --- وهي
محارضة عبثية لكابل بحري؛ وقد نجحت وشائع التحميل في خطوط
الهاتف البرية، فقيمتا $G$ و $R\Gamma$ فيها مختلفتان.
\end{solution}

\begin{solution}{exo:b2:dispersion-wave-packets:10}
(أ) $\partial_te = \dot y(\mu\ddot y + Ky) + Ty'\dot y' = \dot yTy'' + Ty'\dot y' = \partial_x
(Ty'\dot y) = -\partial_x\mathcal P$. (ب) $\langle e\rangle = \tfrac14A^2(\mu\omega^2 + Tk^2 + K) =
\tfrac12\mu\omega^2A^2$ (باستعمال $\mu\omega^2 = Tk^2 + K$)؛ و $\langle\mathcal P\rangle = \tfrac12Tk\omega
A^2$. (ج) والنسبة $Tk/\mu\omega = c^2k/\omega = v_g$. (د) $y = A\eu^{-\kappa x}\cos\omega t$:
ومنه $\mathcal P \propto \sin2\omega t$، ومتوسطه معدوم --- فلا طاقة تعبر الوسط.
\end{solution}

\begin{solution}{exo:b2:dispersion-wave-packets:11}
(أ) $\tanh kh \to 1$: ومنه $\omega^2 = gk$؛ وعند $kh \ll 1$: $\omega = \sqrt{gh}\,k$، والسرعة $\sqrt{gh}$،
وغير متبدّدة. (ب) $kh = 0.13$: أي ضحل؛ و $c = \sqrt{9.81 \times 4000} = \qty{198}{m/s}$،
و $T = \qty{17}{min}$؛ و \qty{8000}{km} في \qty{11}{h}. (ج) $\tfrac12\rho gA^2 = \qty{1.2}{kJ/m^2}$
على $200 \times 1000$ km$^2$: أي \qty{2.5e14}{J}. وبغرين: $A \propto h^{-1/4}$، و $(4000/10)^{1/4}
= 4.5$: أي \qty{2.2}{m}، قبل الانحدار والصعود على الشاطئ. (د) عند $h \approx \lambda/2 =
\qty{50}{m}$؛ ويتباطأ الجزء من الذروة الواقع في ماء أضحل، فتدور
الذروة موازيةً لخطوط تساوي العمق --- أي انكسار.
\end{solution}

\begin{solution}{exo:b2:dispersion-wave-packets:12}
(أ) عند $x = \ell$: $v = v_i + v_r$، و $i = (v_i - v_r)/Z_c$، و $v = R_Li$: ومنه $r = (R_L -
Z_c)/(R_L + Z_c)$. (ب) $0$ و $-1$ و $+1$ و $0.5$. (ج) ويبدو مدخل الخط مثل
$Z_c$، فيرى الراسم \qty{0.5}{V} عند $t = 0$، ثم الصدى $0.5r$ فولت عند
$2\ell/c = \qty{1}{\micro s}$: أي لا شيء، و \qty{-0.5}{V}، و \qty{+0.5}{V}، و \qty{+0.25}{V}؛ ويمتصّ
المولّد المواءم الصدى، فلا يكون هناك صدى ثانٍ. (د)
$\tfrac12 \times 2 \times 10^8 \times 1.8 \times 10^{-6} = \qty{180}{m}$.
\end{solution}

\begin{solution}{pb:b2:dispersion-wave-packets:1}
\textbf{1.} $\partial_xv = -\Lambda\partial_ti - Ri$، و $\partial_xi = -\Gamma\partial_tv$؛ ومع
$\eu^{\iu(\omega t - kx)}$: $\iu kv = (\iu\omega\Lambda + R)i$، و $\iu ki = \iu\omega\Gamma v$، ومنه
$k^2 = \omega\Gamma(\omega\Lambda - \iu R)$.

\textbf{2.} $\omega = R/\Lambda = \qty{4000}{rad/s}$، أي \qty{640}{Hz}: فعند بضعة هرتزات يكون $\Lambda\omega
\ll R$، أي النظام الانتشاري.

\textbf{3.} $k' = k'' = \sqrt{R\Gamma\omega/2}$ مع $R\Gamma = \qty{5e-13}{s/m^2}$: فعند \qty{1}{Hz}:
$\qty{1.25e-6}{m^{-1}}$، و $v_\varphi = \qty{5e6}{m/s}$، و $\delta = \qty{800}{km}$؛ وعند \qty{4}{Hz}:
$\qty{2.5e-6}{m^{-1}}$، و $\qty{1e7}{m/s}$، و \qty{400}{km}.

\textbf{4.} $8.7L/\delta$: أي \qty{38}{dB} عند \qty{1}{Hz}، و \qty{76}{dB} عند \qty{4}{Hz}: فلا تنجو إلا
أبطأ المركّبات؛ وتصل النقطة الحادّة حدبةً طويلة منخفضة.

\textbf{5.} $R\Gamma L^2 = 5 \times 10^{-13} \times 1.2 \times 10^{13} = \qty{6}{s}$: أي نقطة كل
ست ثوانٍ، وكلمة في الدقيقة (وقد بلغ الكابل الحقيقي، بمقاييس
غلفانية مرآوية حساسة، بضع كلمات).

\textbf{6.} $\Lambda = R/\omega = \qty{80}{mH/km}$. وعندئذ $k^2 = \Gamma R\omega(1 - \iu)$، و $k = \sqrt{\Gamma R\omega}
\,2^{1/4}\eu^{-\iu\pi/8}$: ومنه $k' = \qty{3.9e-6}{m^{-1}}$، و $k'' = \qty{1.6e-6}{m^{-1}}$: و $v_\varphi =
\qty{6.4e6}{m/s}$، و $\delta = \qty{620}{km}$، و \qty{49}{dB} على الكابل بدل
\qty{76}{dB}.

\textbf{7.} نقطة كل \qty{6}{s} هي نحو \qty{0.2}{bit/s}: فالليف أسرع
$10^{10}$ إلى $10^{11}$ مرة.

\textbf{8.} $v_\varphi = gT/2\pi$، و $v_g = gT/4\pi$.

\textbf{9.} عند \qty{18}{s}: $v_g = \qty{14.1}{m/s}$، و $4 \times 10^6/14.1 = \qty{3.3}{d}$؛
وعند \qty{6}{s}: \qty{4.7}{m/s}، أي تسعة أيام وتسعة أعشار: فيدوم الموج الطويل أكثر من ستة أيام.

\textbf{10.} $t = D/v_g = 4\pi D/gT$، ومنه $1/T = gt/4\pi D$. و $\Delta(1/T) = 1/12 - 1/16
= \qty{0.0208}{Hz}$ في \qty{86400}{s}: ومنه $g/4\pi D = \qty{2.4e-7}{s^{-2}}$، و $D = \qty{3200}{km}$.

\textbf{11.} $\tfrac12\rho gA^2 = \qty{11}{kJ/m^2}$؛ وتتحرك الطاقة بالسرعة $v_g = gT/4\pi =
\qty{9.4}{m/s}$: أي \qty{103}{kW} في المتر من الذروة.

\textbf{12.} $10^3 \times 103 \times 0.3 = \qty{31}{MW}$: أي عشر عنفات ريحية.

\textbf{13.} $\lambda = 2\pi U^2/g = \qty{16}{m}$. وتنتقل طاقة كل مركّبة موجية
بنصف سرعة طورها، ومنه ينحصر النقش في إسفين
خلف الزورق (نصف زاويته \qty{19.5}{\degree}، مهما تكن سرعته).

\textbf{14.} تولد الذروة في مؤخرة المجموعة، وتتجاوزها
(فالذرى تمضي بالسرعة $2v_g$ بالنسبة إلى الماء، و $v_g$ بالنسبة إلى
المجموعة) وتموت في مقدمتها. وطول المجموعة $120 \times 9.4 = \qty{1.1}{km}$،
أي خمسة أطوال موجة ($\lambda = gT^2/2\pi = \qty{225}{m}$)؛ وتمرّ عشر ذرى بالعوّامة،
وقد عاشت كل واحدة نحو دقيقتين.

\textbf{15.} $kh = 0.13$: أي ضحل. و $c = \qty{198}{m/s}$، و $T = \lambda/c = \qty{17}{min}$،
و \qty{8.4}{h}.

\textbf{16.} غير متبدّد: فكل المركّبات عند \qty{198}{m/s}، وتصل
معًا نبضةً واحدة (وأما الموج الطويل، عند $\sim\qty{10}{m/s}$، مفروزًا بالدور،
فيستغرق أسبوعًا).

\textbf{17.} $\tfrac12\rho gA^2 = \qty{1.8}{kJ/m^2}$ على $500 \times 200$ km$^2$: أي \qty{1.8e14}{J}،
و $4\%$ من ميغاطن.

\textbf{18.} $A = 0.6(4000/20)^{1/4} = \qty{2.3}{m}$؛ و $c = \sqrt{gh} = \qty{14}{m/s}$، و $\lambda =
cT = \qty{14}{km}$. والذروة، وهي في ماء أعمق من القاع الذي أمامها،
تمضي أسرع فتلحق به: فتنحدر الجبهة وتصير جدارًا.

\textbf{19.} ارتفاع \qty{0.6}{m} على امتداد ربع ساعة، على ميل
مقداره $10^{-5}$: فلا شيء يُحسّ.

\textbf{20.} السرعة الوسطى، $\sqrt{gh}$ من أجل العمق الوسطي؛ ولأن
$\sqrt h$ مقعّرة، تبطّئ المقاطع الضحلة الموجةَ أكثر مما تسرّعها
العميقة، ويكون متوسط المسار دون $\sqrt{g\bar h}$.

\textbf{21.} $\omega' = \tfrac12\sqrt{g/k}$، و $\omega'' = -\tfrac14\sqrt{g/k^3}$؛ ومن أجل $k = \qty{0.063}{rad/m}$:
$\omega'' = \qty{-50}{m^2/s}$.

\textbf{22.} $t_{\text{تمدد}} = 10^6/50 = \qty{2e4}{s}$، أي \qty{5.6}{h}؛ وقد تحركت
الرزمة $v_gt = 6.25 \times 2 \times 10^4 = \qty{125}{km}$.

\textbf{23.} الرتل الطويل يحوي حزمة ضيّقة من $k$، وسرعات
مجموعتها متقاربة: فتبقى المركّبات معًا مدةً أطول. ووسط
التسونامي غير متبدّد: فالسرعات متساوية، ولا تمدد.

\textbf{24.} $\omega'' = \hbar/m = \qty{1.2e-4}{m^2/s}$: ومنه $t_{\text{تمدد}} = 10^{-18}/1.2 \times
10^{-4} = \qty{9e-15}{s}$؛ و $v_g = \hbar k/m = \qty{7.3e6}{m/s}$: أي \qty{60}{nm}.

\textbf{25.} الكابل: $k^2 = -\iu R\Gamma\omega$، و $v_\varphi = \sqrt{2\omega/R\Gamma}$، ولا
$v_g$ نظيفة، ويوهن ويلطّخ. والموج الطويل: $\omega^2 = gk$، و $v_\varphi = 2v_g$،
ويتمدد. والتسونامي: $\omega = \sqrt{gh}\,k$، و $v_\varphi = v_g$، ولا هذا ولا ذاك. والإلكترون:
$\omega = \hbar k^2/2m$، و $v_g = 2v_\varphi$، ويتمدد.
\end{solution}
